\documentclass[10pt,a4paper]{amsart}
\usepackage[margin=19mm]{geometry}
\usepackage{amsmath,amssymb,mathtools}
\usepackage[scr=rsfs,cal=euler]{mathalfa}
\usepackage{xfrac}
\usepackage[unicode,hidelinks,bookmarks=false]{hyperref}
\newcommand{\ie}{i.e.\ }
\newcommand{\cf}{cf.\ }
\newcommand{\resp}{respectively\ }
\newcommand{\Subsection}{\subsection}
\providecommand{\nobreakdash}{\nobreak\hbox{-}}
\setlength{\emergencystretch}{2em}
\multlinegap0pt
\usepackage{array}
\usepackage{euscript}
\usepackage{graphics}
\usepackage{cancel}
\usepackage{fancybox}
\usepackage{verbatim}
\usepackage{paralist}
\usepackage{setspace}
\usepackage{fancyhdr}
\usepackage{url}
\usepackage{ulem}
\usepackage{tikz}
\usetikzlibrary{arrows,matrix, patterns}
\renewcommand{\headrulewidth}{0pt}
\newtheorem{thm}{Theorem}[section]
\newtheorem{cor}[thm]{Corollary}
\newtheorem{lem}[thm]{Lemma}
\newtheorem{dff}[thm]{Definition}
\newtheorem{xmp}[thm]{Example}
\newtheorem{exer}[thm]{Exercise}
\newtheorem{rmk}[thm]{Remark}
\newtheorem{conj}[thm]{Conjecture}
\newcommand{\FF}{\mathbf F}
\newcommand{\ZZ}{\mathbf Z}
\newcommand{\QQ}{\mathbf Q}
\newcommand{\NN}{\mathbf N}
\newcommand{\qedbox}{\rule{2mm}{2mm}}
\renewcommand{\qedsymbol}{\qedbox}
\renewcommand{\hat}{\widehat}
\newcommand{\id}{\operatorname{id}}
\newcommand{\fG}{\mathfrak{G}}
\newcommand{\fm}{\mathfrak{m}}
\newcommand{\cA}{\mathcal{A}}
\newcommand{\cE}{\mathcal{E}}
\newcommand{\cF}{\mathcal{F}}
\newcommand{\cI}{\mathcal{I}}
\newcommand{\cJ}{\mathcal{J}}
\newcommand{\cP}{\mathcal{P}}
\renewcommand{\d}{\delta}
\def\ur{\textrm{ur}}
\def\alg{\textrm{alg}}
\def\ux{\underline{x}}
\def\cs{s}
\def\cN{\mathcal{N}}
\def\Gal{\operatorname{Gal}}
\colorlet{DG}{green!50!black}
\colorlet{DO}{orange!50!black}
\colorlet{DR}{red!50!black}
\colorlet{DB}{blue!50!black}
\colorlet{DP}{purple!50!black}
\begin{document}
\title[Arithmetic partial differential operators on ramified extens]{Arithmetic partial differential operators on ramified extensions of $\ZZ_p$}
\author{Marshall Donn, Lance Edward Miller, William D. Taylor}
\date{}
\maketitle
\noindent\emph{Source and licence.} Arithmetic partial differential operators on ramified extensions of $\ZZ_p$. Version arXiv:2608.19612v1 (CC BY 4.0). This material is distributed under the Creative Commons Attribution 4.0 International licence, http://creativecommons.org/licenses/by/4.0/, as recorded in the arXiv OAI licence metadata for this exact version and shown on the abstract page https://arxiv.org/abs/2608.19612v1. Full rights notice: licenses/arxiv-2608.19612-CC-BY-4.0.txt.\par\smallskip
\noindent\textit{Editorial scope.} Counted unit: the original \texttt{thm} environment \texttt{lma1} at source lines 328--335 together with its complete proof at lines 336--364; the delivered document keeps source lines 261--365 so that every referenced label and every citation resolves inside it. Native editable source is the paper's own concatenated TeX; the AMS wrapper is editorial formatting only. No publisher class, upstream script or unrelated artwork is redistributed.
\section{Arithmetic differential operators are analytic} Throughout, we keep a fixed partial $\delta_{\pi}$-structure on $A_\pi$ coming from a set $\{ \phi_1,\ldots,\phi_n\}$ of Frobenius automorphisms. The main result \cite[Thm. 4, 5]{BRS11} establishes that a function $f \colon \ZZ_p \to \ZZ_p$ is an arithmetic differential operation of order at most $r$ if and only if it analytic of level $r$. The by far easier direction of the two is to show any arithmetic differential operation of level $r$ is analytic. In the language here, this calculates $\d_\mu(a+\pi^r u)$ as a power series expression, see \cite[Lem. 10]{BRS11} from which it is easy to ascertain that this power series is $\pi$-adically restricted. 

\

To generalize this direction to the PDE setting, we need the following lemma, which is an enhancement of \cite[Lem. 10]{BRS11}. However, the combinatorics needed to accurately describe $\d_\mu(a + \pi^r u)$ as a power series in the PDE setting are significantly more complicated, so we treat the lemma after developing some helpful notation. 

\

\begin{xmp}\label{xmp:n=2} It is helpful to see the complication in the case $n = 2$. A direct computation shows the shape of $\d_2(\d_1(a + \pi^r u))$. First we have from definition
\begin{eqnarray*}
\d_1( a + \pi^r u) & = & \frac{1}{\pi} \left( \phi_1(a + \pi^r u) - (a + \pi^r u)^p  \right) \\
& = & \d_1(a) + \sum_j c_j u^j + \frac{\phi_1(\pi)^r}{\pi} \phi_1(u). 
\end{eqnarray*} In this expression, each $c_j$ is divisible by $\pi^r$. As $\phi_1(\pi) \in (\pi)$ we may write this as $$\d_1(a+\pi^r u) = \d_1(a) + \pi^{r-1}  \sum_{i+j \geq 1} c(i,j) u^i \phi_1(u)^j, \textrm{ for }c(i,j) \in A_\pi.$$ 
Furthermore, expanding once again, we find an even more unwieldy expression 
{\footnotesize
\begin{eqnarray*}
\d_2\left( \d_1(a) + \pi^{r-1}  \sum c(i,j) u^i \phi_1(u)^j \right) & = & \frac{1}{\pi} \left( \phi_2\left( \d_1(a) + \pi^{r-1}  \sum c(i,j) u^i \phi_1(u)^j \right) - \left( \d_1(a) + \pi^{r-1}  \sum c(i,j) u^i \phi_1(u)^j \right)^p  \right) \\
& = &  \d_2(\d_1(a)) + \pi^{r-2}\sum_{i,j,k,\ell} c(i,j,k,\ell) u^i \phi_1(u)^j \phi_2(u)^k \phi_2(\phi_1(u))^\ell
\end{eqnarray*}} for values $c(i,j,k,\ell) \in A_\pi$. 

\end{xmp}


\

The point that Example~\ref{xmp:n=2} shows is that the expression $\d_\mu(a + \pi^m u)$ is not going to be a power series in $u$ alone, since the functions $\phi_i$ will not always be identity as was the case in \cite{BRS11}. Instead, these are complicated power series in $\phi_\nu(u)$ for all subwords $\nu \leq \mu$ and we want to introduce a more convenient notation to keep track of the terms. 

\

The key to simplifying the notation is to write these as sums over all functions from one set to another. As usual, for sets $A$ and $B$, the set of functions from $A$ to $B$ can be denoted $B^A$. For a function $\beta\in B^A$, if $B$ is a subset of the real numbers, then we set $|\beta|_1:=\sum_{a\in A}|\beta(a)|$ if $\beta(a)$ is nonzero for only finitely many inputs and $\infty$ otherwise.

\

\begin{dff}
For each $\mu \in \mathbb{M}^r$, we define the following. 
\begin{enumerate}
\item  Let $\mathcal{P}_\mu := \{ \nu \in \mathbb{M}^r \mid \nu \leq \mu \}$ be the set of subwords of $\mu$, including of course the empty subword. 
\item Set $j_{\text{const}} \colon \mathcal{P}_\mu \to \mathbf{N}$ the zero function. 
\item For $\nu\in\mathcal{P}_\mu$, the indicator function $j_\nu \colon \mathcal{P}_\mu \to \mathbf{N}$ is the function given by $j_\nu(\sigma)=0\text{ for }\sigma\neq \nu$ and $j_\nu(\nu) = 1$. 
\item Let $$J_\mu := \left\{j \colon \mathcal{P}_\mu\to \mathbf{N}\mid j(\mu)=0\text{ and } |j|_1\leq p^{|\mu|}\right\} \cup \{j_\mu\}$$ and note $j_{\text{const}} \in J_\mu$. 
\end{enumerate}
\end{dff} 

Thus, repeating and expanding the calculation from Example~\ref{xmp:n=2}, we see for each $\mu \in \mathbb{M}^r$ and $a \in A_\pi$ there exist $c_a(j,\mu)\in A_\pi$ such that \[\delta_\mu(a + \pi^r u) = \sum_{j\in J_\mu}c_a(j,\mu)\prod_{\nu\in \mathcal{P}_\mu}\phi_\nu(u)^{j(\nu)}.\] 

 \ 
Immediately, we have a relationship between arithmetic partial differential operators and analytic functions.

\begin{thm}\label{thm:PADOtoAnal} Every arithmetic partial differential operator of order at most $r$ is also analytic of level $r$. Specifically, set $\cA$ a collection of centers of balls of radius $1/\pi^r$ covering $A_\pi$ and $F\in \Gamma$. Set also $\tau_a \colon \Gamma\to \Gamma$ the $A_\pi$-algebra homomorphism induced by $\tau_a(x_\mu) := \sum_{j \in J_\mu} c_a(j,\mu) \prod_{\nu \in \mathcal{P}_\mu} x_\nu^{j(\nu)}$. We have that $\Delta(F) = \Phi^{\oplus \cA}(\tau_a(F))_{a\in \cA}$. 
\end{thm}

\begin{proof}
 Write $\delta_\mu(a+\pi^r u) = \Phi(\tau_a(x_\mu))(u)$.  Therefore, for any $a, u\in A_\pi$, 
\[	\Delta(F)(a+\pi^r u) = \Phi(\tau_a(F))(u).\qedhere\]
\end{proof}


Similar to \cite[Lem. 10]{BRS11} we estimate the $\pi$-adic absolute values of $c(j,\mu)$. 


\







\begin{thm}\label{lma1} Fix $r \geq 0$. If $a \in A_{\pi}$, $v = a + \pi^r u$ in the disc $a+ \pi^r A_{\pi}$, and $\mu \in \mathbb{M}^r$, expanding, then $$\delta_{\mu}(v) = \sum_{j \in J_\mu} c_a(j,\mu) \prod_{\nu \in \mathcal{P}_\mu} \phi_\nu(u)^{j(\nu)}$$ with $c_a(j,\mu) \in A_\pi$ satisfying the following:
\begin{enumerate}
\item $|c_a(j_{\textup{const}},\mu)|_\pi \leq 1$, and in particular $c_a(j_\textup{const},\mu) = \delta_\mu(a)$,
\item $|c_a(j_{\nu},\mu)|_\pi \leq \frac{1}{\pi^{r-|\nu|}}$ for all $\nu\in\mathcal{P}_\mu$, with equality when $\nu = \mu$,
\item $|c_a(j,\mu)|_\pi \leq \frac{1}{\pi^{(r-|\mu|+1)|j|_1-1}}$ when $|j|_1\geq 2$,
\end{enumerate}
In particular, if $j\neq j_\textup{const}$, then $c_a(j,\mu)$ is a unit if and only if $j = j_\mu$ and $|\mu|= r$.
\end{thm}

\begin{proof} As $a$ is fixed throughout, we suppress this subscript from the notation. We proceed by induction on $|\mu|$. The base case $|\mu| = 0$ is immediate. Thus, we proceed with the inductive case. Assume that the estimates hold for $\mu$. Fix $1 \leq i \leq n$. We show the estimates hold for $i \cdot \mu$. By definition

\begin{align*}
	\delta_{i\cdot \mu}(a + \pi^ru)
	&= \dfrac{1}{\pi}\left(\phi_i\left(\sum_{j\in J_\mu} c(j,\mu)\prod_{\nu\in P_\mu} \phi_\nu(u)^{j(\nu)}\right) - \left(\sum_{j\in J_\mu} c(j,\mu)\prod_{\nu\in P_\mu} \phi_\nu(u)^{j(\nu)}\right)^p\right)\\
	&=\dfrac{1}{\pi}\left(\phi_i(c(j_{\mathrm{const}},\mu)) + \phi_i(c(j_\mu,\mu))\phi_{i\cdot\mu}(u) + \sum_{j\in J_\mu\colon j\neq j_{\mu,j_{\mathrm{const}}}}\phi_i(c(j,\mu))\prod_{\nu\in P_\mu} \phi_{i\cdot\nu}(u)^{j(\nu)} \right.\\
	&\qquad - \left.\left(c(j_{\mathrm{const}},\mu) + c(j_\mu,\mu)\phi_{\mu}(u) + \sum_{j\in J_\mu\colon j\neq j_{\mu,j_{\mathrm{const}}}}c(j,\mu)\prod_{\nu\in P_\mu} \phi_{\nu}(u)^{j(\nu)}\right)^p\right)
\end{align*}

\

Statement (1) and the $\nu = \mu$ case of statement (2) follow by direct computation as 
\[c(j_{\textup{const}},i \cdot \mu) = \dfrac{\phi_i(c(j_{\mathrm{const}},\mu)) -c(j_{\mathrm{const}},\mu)^p}{\pi} = \d_i(c(j_{\mathrm{const}},\mu)) = \d_i(\d_\mu(a)) = \d_{i \cdot \mu}(a)\] and 
$$c(j_{i \cdot \mu},i \cdot \mu ) = \frac{1}{\pi}\phi_i(c(j_\mu,\mu)) =  \frac{1}{\pi} \phi_i(\phi_\mu(\pi^r)/\pi^{|\mu|})$$ which forces $|c(j_{i \cdot \mu},i \cdot \mu )|_\pi = \left| \frac{1}{\pi^{|\mu|+1}} \phi_{i \cdot \mu}(\pi^r) \right|_\pi = \frac{1}{\pi^{r - |\mu| - 1}} =\frac{1}{\pi^{r - |i \cdot \mu|}}$.
 
\

For the remainder of (2), let $\nu\in\mathcal{P}_{i\cdot \mu}\setminus \{i\cdot \mu\}$.  If $\nu =i\cdot  \nu'$ and $\nu'\in \mathcal{P}_\mu$, then $c(j_\nu,i\cdot \mu)$ has as a summand $\frac{1}{\pi}\phi_i(c(j_{\nu'}, \mu))$.  By the inductive hypothesis, 
\[\left|\frac{1}{\pi}\phi_i(c(j_{\nu'}, \mu))\right|_\pi  = \pi\left|c(j_{\nu'}, \mu)\right|_\pi\leq \frac{1}{\pi^{r-|\nu'| - 1}} = \frac{1}{\pi^{r-|\nu|}}.\]
If $\nu\in \mathcal{P}_\mu$, then $c(j_\nu,i\cdot \mu)$ has as a summand $\frac{p}{\pi}c(j_\textup{const})^{p-1}c(j_\nu,\mu)$, and by the inductive hypothesis,
\[\left|\frac{p}{\pi}c(j_\textup{const})^{p-1}c(j_\nu,\mu)\right|_\pi \leq |c(j_\nu,\mu)|_\pi \leq \frac{1}{\pi^{r-|\nu|}}.\]

Now, for assertion (3), suppose that $j\in J_{i \cdot \mu}$ and $|j|_1\geq 2$. All summands of $c(j,i \cdot \mu)$ have only two forms. One form is \begin{equation*}
\frac{c(j^{(1)},\mu) \cdots c(j^{(p)},\mu)}{\pi}\end{equation*}
for some $j^{(1)},\ldots,j^{(p)} \in J_\mu$ with $\sum_{k=1}^p |j^{(k)}|_1 = |j|_1$. Note,
 we have no control over what functions $j^{(k)}$ can be. Yet, in all cases, using the inductive hypothesis we see $|c(j^{(k)},\mu)|_\pi \leq \frac{1}{\pi^{(r-|i\cdot \mu| + 1)|j^{(k)}|_1}}$ for each $k$. Therefore,  
\[\left|\frac{c(j^{(1)},\mu) \cdots c(j^{(p)},\mu)}{\pi}\right|_\pi =\pi\prod_{k=1}^p|c(j^{(k)},\mu)|_\pi\leq \frac{\pi}{\pi^{\sum_{k=1}^p(r-|i\cdot \mu|+1)|j^{(k)}|_1}} = \frac{1}{\pi^{(r-|i\cdot \mu|+1)|j|_1-1}}.\] The other type of summand only occurs when all of the words in the support of $j$ are of the form $i\cdot \nu$. In this case, $c(j,i\cdot \mu)$ also has summands of the form $\frac{\phi_i(c(\tilde j,\mu))}{\pi}$, where $\tilde j\in J_\mu$ is defined by $\tilde j(\nu) = j(i\cdot \nu)$.  In particular, this means that $|\tilde j|_1 = |j|_1\geq 2$.  Therefore by the inductive hypothesis,
\[\left|\frac{\phi_i(c(\tilde j,\mu))}{\pi}\right|_\pi \leq \pi\left|c(\tilde j,\mu)\right|_\pi \leq \frac{1}{\pi^{(r-|\mu|+1)|\tilde j|_1 - 2}}\leq  \frac{1}{\pi^{(r-|i\cdot \mu| + 1)|j|_1 + |j|_1 -2}}\leq \frac{1}{\pi^{(r-|i\cdot \mu| + 1)|j|_1 - 1}}.\qedhere\]
\end{proof}
    

\begin{thebibliography}{99}
\bibitem[BRS11]{BRS11} A. Buium, C. Ralph, and S. Simanca , {\em Arithmetic differential operators on $\ZZ_p$}, Journal of Number Theory, 131, (2011), 96--105.
\end{thebibliography}
\end{document}
